% Ported verbatim from paper/top3_joint_moment_bound.tex lines 957-1084.
% Only cross-reference wording was adapted; no mathematics was altered.
Throughout this section assume $0<r_-\le R\le r_+<\infty$ almost surely,
with $r_+/r_-\le3$, and work in the branch $D<1/20$.
Equation~\eqref{budget4} implies
\begin{equation}\label{smallbranch}
 e<1/20,\quad x>9/10,\quad V<21/104<2/9,\quad C>5/3,
 \qquad q^2=\frac{3-V}{9(1-e)}\le\frac{20}{57}.
\end{equation}
In particular $\mathbf m\ne0$. Put $p_i=\langle T_i,\Psi_i(v)\rangle
=\E[Rg_i(t)]$. The elementary Hermite identities
\begin{align}\label{Hermite}
 \left(\frac{53}{28}t^2-\frac94t+\frac{41}{70}\right)-g_4(t)
       &=(1-t)(t-1/2)^2(t+2)\ge0,\notag\\
 (2t^2-37t/20+1/4)-g_3(t)
       &=(1-t)(t-1/2)^2\ge0
\end{align}
hold on $[-1,1]$. Define
\[
 \gamma_4=9/4-53\sqrt3/42,\quad
 \ell_4=9/4-511\sqrt3/420,\quad
 \ell_3=37/20-11\sqrt3/12.
\]
Here $0<\gamma_4<1/15$ (square the positive rational comparisons).
Since $\E Rt=x$, $\E Rt^2=L+2x/\sqrt3-C/3$, the first identity gives
\[
 -p_4\ge\gamma_4x+19C/420-53L/28
       \ge\ell_4-e/15-19V/1400-53L/28.
\]
We used $x\ge1-e$ and $C\ge\sqrt3-3V/10$; the latter follows from
$V=(\sqrt3-C)(\sqrt3+C)$ and $C>5/3$.
The second identity gives
\[
 -p_3\ge(37/20-4/\sqrt3)x+5C/12-2L
       \ge\ell_3+9e/40-V/8-2L.
\]
Here $4/\sqrt3-37/20>9/20$ and $x\le1-e/2$.
Define the nonnegative differences in these last two inequalities by
\begin{align*}
 r_4&=-p_4-\ell_4+e/15+19V/1400+53L/28,\\
 r_3&=-p_3-\ell_3-9e/40+V/8+2L.
\end{align*}
The polynomial identities, rather than a fixed sphere assumption,
justify these directional estimates for the actual weighted moments.

Set $\eta=1/10$, $\theta=6/25$, and
$\ell=53\eta/28+2\theta=937/1400$. Define the tensor squares
\[
 {\cal R}_4=a_4\left\|T_4+\frac{\eta}{2a_4}\Psi_4(v)\right\|^2,\qquad
 {\cal R}_3=a_3\left\|T_3+\frac{\theta}{2a_3}\Psi_3(v)\right\|^2.
\]
Their exact expansions give
\[
 a_4K^2=-\eta p_4-\frac{2\eta^2}{35a_4}+{\cal R}_4,\qquad
 a_3H^2=-\theta p_3-\frac{\theta^2}{10a_3}+{\cal R}_3.
\]
Define
\[
 E=\frac{5183683}{5659500},\quad
 V_c=\frac{145858703}{581042000},\quad
 J_c=\frac{7989533}{14526050},\quad
 {\cal R}_0={\cal R}_3+{\cal R}_4+\eta r_4+\theta r_3\ge0.
\]
Substituting the directional identities and \eqref{payment} into
\eqref{budget7}, then writing $L=\overline L-r_{\rm reg}$, gives the exact identity
\begin{equation}\label{certificate-start}
 D=d_0+E e+V_c V+J_c J-\ell\overline L
             +{\cal R}_0+r_{\rm rad}+\ell r_{\rm reg}.
\end{equation}
Here
\[
 d_0=\delta_{\mathrm q}+\eta\ell_4+\theta\ell_3
                 -\frac{2\eta^2}{35a_4}-\frac{\theta^2}{10a_3},
\]
which simplifies to \eqref{constants}. The coefficient checks are
\[
 E=b_e-\eta/15+9\theta/40,\quad
 V_c=b_V-19\eta/1400-\theta/8+3/50,\quad J_c=b_J-27/50.
\]
All the remainders in this identity refer to the same probability law.

Put $B_c=V_c-5\ell/18$ and $C_c=J_c-\ell/4$. The exact rational values are
\[
 B_c=\frac{42565069}{653672250},\qquad
 C_c=\frac{88944217}{232416800}.
\]
For $u=\sqrt e$, $z=\sqrt V$, and $j=\sqrt J$, the remaining quadratic has
the explicit completion
\begin{align}\label{squares}
 E e+V_c V+J_c J-\ell\overline L
 &=\frac{13}{200}\left(z-\frac{937}{273}qu\right)^2
   +\frac{19}{50}\left(j-\frac{2811}{2128}qu\right)^2\notag\\
 &\quad+\left(\frac{183}{200}-\frac{725029049}{278853120}q^2\right)e\notag\\
 &\quad+\frac{10481}{11319000}e
        +\frac{305491}{2614689000}V+\frac{625833}{232416800}J.
\end{align}
Indeed the last three fractions are $E-183/200$, $B_c-13/200$,
and $C_c-19/50$. Also
\[
 \frac{725029049}{278853120}
  =\frac{7\ell}{4}+\frac{\ell^2}{9(13/200)}
                          +\frac{9\ell^2}{16(19/50)}.
\]
By \eqref{smallbranch}, the coefficient of $e$ in parentheses is at least
\[
 \frac{183}{200}-\frac{20}{57}\frac{725029049}{278853120}
                         =\frac{53754367}{19868284800}>0.
\]
Thus every term on the right of \eqref{certificate-start} after $d_0$
is nonnegative. The displayed squares give an explicit certificate for the remaining
quadratic form.

\begin{proposition}\label{conditional-law}
Every isotropic law on $\R^3$ whose radii lie almost surely in a positive
interval of endpoint ratio at most $3$ satisfies $D\ge d_0$.
In the branch $D<1/20$, \eqref{certificate-start}--\eqref{squares}
give the full common-remainder certificate.
\end{proposition}
\begin{proof}
The small branch was proved above. Otherwise $D\ge1/20>d_0$.
For explicit exact comparisons use
$682/305<\sqrt5<161/72$ and $433/250<\sqrt3<97/56$,
each verified by squaring. They yield
\[
 d_0-7/144>\frac{940735364069}{29456638041600000}>0,\qquad
 1/20-d_0>\frac{620701786231}{482895705600000}>0.
\]
No $B=0$ division occurs: it is excluded by $e<1/20$ in the small branch.
\end{proof}
